\clearpage
\section{未采用的后续候选}\label{sec:explorations}
BI-A50 至 BI-A58 探索了厂商分页接口、decode 分段配置和组件组合。表~\ref{tab:bi-late} 按验证层级列出结果；接口探针、微基准和端到端对照的证据强度不同。

\begin{table}[!htbp]
\centering\small
\caption{BI-A50--BI-A58：优化内容、验证层级与处置（均未纳入最终实现）\label{tab:bi-late}}
\begin{tabularx}{\linewidth}{@{}l Y Y >{\raggedright\arraybackslash}p{2.0cm}@{}}
\toprule
编号 & 尝试的优化 & 验证层级与结果 & 处置 \\
\midrule
BI-A50 & 厂商 \path{ixformer.batch_paged_attention} 替换 prefill 的 K/V 双路分页收集，省掉 KV 连续化拷贝 & 接口探针：返回无详情 RuntimeError，无有效性能数据 & 接口阻塞 \\
BI-A51 & 厂商 \path{vllm_paged_attention_mtp}（多查询 / GQA paged prefill） & 接口探针：要求 stride$(-2)=$head\_size，当前 NHD 页缓存为 1，调用前被拒 & 布局不兼容 \\
BI-A52 & 页对齐 decode 内核 NSEG 由 32 改 16 & 微基准：等长 4k 快 1.47--1.57\%，16k 回退 2.3--2.7\% & 否决 \\
BI-A53 & 每序列 KV 分段数 NSEG 取 4/8/16/32 的扫描 & 微基准：16 最佳，仅快 1.4--1.6\%，更小分段并行度不足 & 否决 \\
BI-A54 & 混合批 metadata 索引缓存（A45）与 A48 组合 & 官方对照：4k $-0.4745\%$、16k $+0.0308\%$，平均吞吐 $-0.2397\%$ & 不合入 \\
BI-A55 & LibEntry 与连续 copy 快速路径（A46/A47）与 A48 组合 & 官方 16k 对照：$-0.478\%$ & 不合入 \\
BI-A56 & 审计 Inductor 编译产物中的融合机会 & 仅图审计：RoPE、residual+RMSNorm 与 A48 前置链已融合；未跑 benchmark & 关闭 \\
BI-A57 & 默认日志级别 DEBUG 改 INFO，削减 CPU 侧日志开销 & 官方 4k 中位数：$-0.2394\%$ & 不合入 \\
BI-A58 & decode 双累加器 GEMM 的 warps=16 调优 & 微基准快 3.64\%；官方 4k 中位数 $-0.0336\%$ & 不合入 \\
\bottomrule
\end{tabularx}
\end{table}

A50/A51 未取得有效性能结果，不能视作性能回退。A52/A53 只有微基准数据，A56 只有编译图审计；其余候选也只覆盖表中标明的官方场景。上述证据不足以支持采用这些改动。原始记录按实验编号保存在 \path{docs/experiments/runs/}。

\subsection{C500：量化与投机的复核}\label{sec:quant-spec}
赛题允许自行开发的量化策略，但禁止通过修改 \texttt{vllm serve} 超参开启量化或投机采样；
投机采样还需要 MiniCPM5-2B 以外的 draft 模型，与模型限制冲突，因此不作为交付路径。
为避免只凭历史结论排除，本队在最终栈上重测了权重量化的收益上限。
在 C500 上对四个线性形状比较 bf16 与 int8（\texttt{torch.\_int\_mm}）：

\begin{itemize}[leftmargin=2em]
\item M=64（decode）：int8 GEMM 为 33.8/32.8/63.1/58.9 $\mu$s，
bf16 为 33.3/30.0/56.3/53.4 $\mu$s，int8 并不更快；
\item M=2048（prefill）：int8 GEMM 快 1.25--1.42 倍，
但 per-token 激活量化约 165 $\mu$s、int32 反量化为 76--358 $\mu$s，
两步合计使总时间反而变为 bf16 的 1.7--2.6 倍。
\end{itemize}

decode 受权重与 KV 带宽约束，prefill 受 attention 与反量化开销约束，在这套栈上
W8A8 没有可兑现的端到端空间；BI-V150 上的 A27 也给出同向结论（隔离内 1.82 倍，
接入引擎后 $-19.1\%$）。据此关闭量化方向，不再进入候选。微基准脚本与输出见
\path{docs/experiments/runs/c500-quant-probe/}。

\subsection{C500：A48 门控融合的移植评估}\label{sec:a48-port}
egin{sloppypar}
BI-V150 的 A48 用双 FP32 累加器把 gate/up GEMM 与 SwiGLU 融在一个 Triton 内核里，
省掉中间激活的写回。这套做法搬到 C500 的前提是 GEMM 也由 Triton 执行，
但 C500 的 gate\_up 走 mcBLAS（152--197 TFLOPS），且 mcBLAS 的 epilogue 只有
RELU/BIAS/GELU/HARDSWISH，没有 SiLU，也没有双输入乘法，无法在库内完成
\(\mathrm{silu}(gate)\cdot up\)。

在 C500 上实测（M=2048、N=12288、K=2048，bf16）：

\begin{itemize}[leftmargin=2em]
\item 现状：mcBLAS GEMM 522 $\mu$s 加独立 \texttt{silu\_and\_mul} 175 $\mu$s
（151 MB 流量，862 GB/s），合计 697 $\mu$s/层；
\item 融合 Triton：同形状 GEMM 为 707 $\mu$s，epilogue 在寄存器内完成，约 707 $\mu$s/层。
\end{itemize}

融合省下的 175 $\mu$s 小于 GEMM 变慢的 185 $\mu$s，净增约 10 $\mu$s/层
（42 层约 $+0.4$ ms/步，$+0.3\%$）。因此不做移植：A48 的收益依赖 Triton GEMM，
C500 的 GEMM 引擎是 mcBLAS，两者的 epilogue 能力不可叠加。
\end{sloppypar}
